\chapter{Pendahuluan}
\section{Latar Belakang}
Tulis latar belakang di sini. Contoh sitasi: \textcite{Sheng_2026} atau \parencite{Sheng_2026}.

\section{Masalah Penelitian}
Studi ini mengusulkan kerangka Decision-Focused Learning (DFL) end-to-end untuk sistem HVAC fixed-frequency. 
Masalah kontrol dirumuskan sebagai Mixed-Integer Program dengan tiga arsitektur prediktif: 
jaringan RC berbasis fisika, Input-Convex Neural Networks (ICNN), dan ANN non-linear. 
Untuk mengatasi non-diferensiabilitas variabel keputusan diskret, digunakan continuous relaxation 
dengan Implicit Differentiation untuk formulasi konveks dan Stochastic Smoothing untuk lanskap non-konveks, 
ditambah strategi fine-tuning untuk mengurangi parameter terlatih. Validasi dalam lingkungan BOPTEST 
menunjukkan varian DFL meningkatkan kinerja closed-loop, dengan tambahan penghematan energi hingga 8,6\% 
dan perbaikan kenyamanan termal hingga 23,5\% dibanding MPC konvensional. 

Kerangka DFL end-to-end menggabungkan lapisan prediksi dan optimasi MPC. Tiga model prediktif 
(3R2C RC, ANN MLP dua hidden layer, ICNN) dilatih dengan loss prediksi (baseline) atau loss tugas kontrol 
(DFL). Implicit Differentiation via CVXPYlayers dipakai untuk RC-MPC yang direlaksasi kontinu; Stochastic 
Smoothing dengan score-function estimator dan Monte Carlo dipakai untuk ANN/ICNN, dengan fine-tuning layer 
output. Evaluasi closed-loop dalam simulator BOPTEST bestest\_hydronic\_heat\_pump, dua skenario cuaca, 
metrik MAE/RMSE/MAPE/CVRMSE dan KPI biaya, discomfort, switching, waktu komputasi.
